% ── Dashboard ─────────────────────────────────────────────────
\chapter{Dashboard}

        Dashboard sme navrhli ako centrálny ovládací panel aplikácie ProMat, ktorý poskytuje
používateľom rýchly prehľad aktuálneho stavu ich projektov, úloh a~tímov. Slúži ako vstupný
bod po prihlásení a~umožňuje rýchly prístup k~kľúčovým metrikám a~aktívnym projektom.

\section{Prehľad a~úvodné informácie}

        Hlavnú stránku dashboarda sme namapovali na endpoint \texttt{/dashboard}. Zobrazuje
komplexný prehľad pracovného priestoru používateľa. V~hornej časti sa zobrazujú tri kľúčové metriky:

\begin{itemize}
  \item \textbf{Aktívne projekty} -- počet aktívnych projektov, do ktorých má používateľ prístup,
  \item \textbf{Blížiace sa termíny} -- počet úloh s~termínom vypršania v~nasledujúcich 7 dňoch,
  \item \textbf{Aktuálny dátum} -- aktuálny dátum pre rýchlu orientáciu.
\end{itemize}

        Tieto informácie sme zobrazili vo forme kariet v~prehľadnom gride s~vizuálnym odlíšením.

\section{AI-asistované poznatky}

        Jedným z~kľúčových prvkov dashboarda sú \textbf{AI Insights} -- automaticky generované
upozornenia a~poznatky na základe aktuálneho stavu úloh a~projektov. Systém dynamicky
vytvára maximálne 4 poznatky naraz podľa priority.

\subsection{Typy poznatkov}

        Medzi hlavné typy poznatkov patria:

\begin{itemize}
  \item \textbf{Detekcia rizík} -- upozornenie na úlohy po termíne,
  \item \textbf{Upozornenie na termín} -- upozornenie na úlohy splatné v~najbližších 24 hodinách,
  \item \textbf{Plánovaná práca} -- prehľad plánovaných úloh v~horizonte nasledujúceho týždňa,
  \item \textbf{Úspech} -- zvýraznenie nedávno dokončených úloh,
  \item \textbf{Vyžaduje pozornosť} -- signalizácia vysoko prioritných úloh bez stanoveného termínu.
\end{itemize}

        Každý poznatkok obsahuje popis, odkaz na podrobnosti a~vizuálne odlíšenie farbou,
ktorá indikuje typ upozornenia (obr.~
ef{fig:dashboard_statistiky}).

\section{Sledovanie prioritných projektov}

        Sekciu sledovania prioritných projektov sme navrhli tak, aby zobrazovala top 3 aktívne projekty
usporiadané podľa času poslednej úpravy. Pri každom projekte sa uvádza názov, percento pokroku,
textové označenie termínu (napr. blíži sa termín alebo po termíne), farebný indikátor stavu
a~počet členov tímu.

        Používateľ môže kliknúť na projekt a~prejsť na stránku s~jeho podrobnosťami.

\section{Vyhľadávanie}

        Do aplikácie sme zaradili globálnu vyhľadávaciu funkciu dostupnú na endpoint \texttt{/search}.
Používateľ môže rýchlo vyhľadávať projekty podľa názvu alebo popisu, úlohy podľa názvu alebo
popisu a~dokumenty podľa názvu alebo obsahu.

        Vyhľadávanie je necitlivé na veľkosť písmen a~aplikuje sa na projekty viditeľné
danému používateľovi (na základe oprávnení a~členstva v~tímoch). Výsledky sú zoradené
chronologicky podľa času vytvorenia alebo aktualizácie.

\section{Bezpečnosť a~oprávnenia}

        Všetky zobrazené údaje na dashboarde sme filtrovali podľa oprávnení používateľa.
Funkcia \texttt{\_visible\_project\_ids\_query()} zabezpečuje, že používateľ vidí iba projekty,
ku ktorým má oprávnený prístup. Systém automaticky aplikuje pravidlá členstva tímu
a~všetky metriky aj poznatky sa počítajú iba z~viditeľných projektov a~úloh (obr.~
ef{fig:dashboard_hlavny}).

\begin{figure}[H]
  \centering
  \includegraphics[width=0.9\textwidth]{images/dashboard.png}
  \caption{Dashboard aplikácie ProMat}
  \label{fig:dashboard_hlavny}
\end{figure}

\begin{figure}[H]
  \centering
  
\includegraphics[width=0.9\textwidth]{images/dashboard_statistiky.png}
  \caption{Štatistiky a~AI-asistované poznatky na dashboarde}
  \label{fig:dashboard_statistiky}
\end{figure}
